\chapter{2015年湖北卷（文）}

\section{单选题}

\begin{problem}
i 为虚数单位，\( i^{607}= \)（\quad）

\choices
    {i}
    {-i}
    {1}
    {\(-1\)}
\end{problem}

\begin{problem}
我国古代数学名著《数书九章》“有米经分数”题: 粮仓开仓收粮，有人送来米 1534 石，装得米内夹谷，抽样取米一把，数得 254 粒内夹谷 28 粒，则这批米内夹谷约为（\quad）

\choices
    {134 石}
    {169 石}
    {338 石}
    {1365 石}
\end{problem}

\begin{problem}
命题“ \( 3\pi_{2}\in(0,+\infty) \) ，\( \ln a_{2}=\pi_{2}-1 \) ”的否定是（\quad）

\choices
    {\(\exists \pi_2\in (0, + \infty),\ln a_2\neq a_2 - 1\)}
    {\(\exists \pi_2\notin (0, + \infty),\ln a_2 = \pi_2 - 1\)}
    {\(\forall x\in (0, + \infty),\ln x\neq x - 1\)}
    {\(\forall x\notin (0, + \infty),\ln x = x - 1\)}
\end{problem}

\begin{problem}
已知变量 \(x\) 和 \(y\) 满足关系 \(y = -0.1x + 1\)，变量 \(y\) 与 \(z\) 正相关. 下列结论中正确的是（\quad）

\choices
    {\(x\) 与 \(y\) 正相关，\(x\) 与 \(z\) 负相关}
    {\(x\) 与 \(y\) 正相关，\(x\) 与 \(z\) 正相关}
    {\(x\) 与 \(y\) 负相关，\(x\) 与 \(z\) 负相关}
    {\(x\) 与 \(y\) 负相关，\(x\) 与 \(z\) 正相关}
\end{problem}

\begin{problem}
\(l_{1}, l_{2}\) 表示空间中的两条直线. 若 \(p: l_{1}, l_{2}\) 是异面直线; 若 \(l_{1}: l_{2}\) 不相交，则（\quad）

\choices
    {\(p\) 是 \(q\) 的充分条件，但不是 \(q\) 的必要条件}
    {\(p\) 是 \(q\) 的必要条件，但不是 \(q\) 的充分条件}
    {p 是 q 的充分必要条件}
    {\(p\) 既不是 \(q\) 的充分条件，也不是 \(q\) 的必要条件}
\end{problem}

\begin{problem}
函数 \( f(x) = \sqrt{4 - |x|} + \lg \frac{x^2 - 5x + 6}{x - 3} \) 的定义域为（\quad）

\choices
    {(2,3)}
    {(2,4)}
    {\((2,3)\cup (3,4)\)}
    {\((-1,3)\cup (3,6]\)}
\end{problem}

\begin{problem}
已知符号函数 \(\operatorname{sgn} x = \left\{ \begin{array}{ll} 1, & x > 0, \\ 0, & x = 0, \\ -1, & x < 0, \end{array} \right.\)（\quad）

\choices
    {\( |x| = x|\operatorname{sgn} x| \)}
    {\( |x| = x\operatorname{sgn}|x| \)}
    {\( |x| = |x|\operatorname{sgn} x \)}
    {\( |x| = x\operatorname{sgn}x \)}
\end{problem}

\begin{problem}
在区间 \([0,1]\) 上随机取两个数 \(x, y\)，记 \(p_1\) 为事件“\(x + y \leqslant \frac{1}{2}\)”的概率，\(p_2\) 为事件“\(xy \leqslant \frac{1}{2}\)”的概率，则（\quad）

\choices
    {\(p_1 < p_2 < \frac{1}{2}\)}
    {\(p_2 < \frac{1}{2} < p_1\)}
    {\(\frac{1}{2} < p_2 < p_1\)}
    {\(p_1 < \frac{1}{2} < p_2\)}
\end{problem}

\begin{problem}
将离心率为 \( e_1 \) 的双曲线 \( C_1 \) 的实半长轴 \( k \) 和虚半轴长 \( b \) （\( a \neq b \)）同时增加 \( m (m > 0) \) 个单位长度，得到离心率为 \( e_2 \) 的双曲线 \( C_2 \)，则（\quad）

\choices
    {对任意的 \(a, b, e_1 > e_2\)}
    {当 \(a > b\) 时，\(e_1 > e_2\); 当 \(a < b\) 时，\(e_1 < e_2\)}
    {对任意的 \(a, b, e_1 < e_2\)}
    {当 \(a > b\) 时，\(e_1 < e_2\); 当 \(a < b\) 时，\(e_1 > e_2\)}
\end{problem}

\begin{problem}
已知集合 \(A = \{(x, y) | x^2 + y^2 \leqslant 1, x, y \in \Z\}\)，\(B = \{(x, y) | x| \leqslant 2, |y| \leqslant 2, x, y \in \Z\}\)，定义集合 \(A \oplus B = \{(x_1 + x_2, y_1 + y_2) | (x_1, y_1) \in A, (x_2, y_2) \in B\}\)，则 \(A \oplus B\) 中元素的个数为（\quad）

\choices
    {77}
    {49}
    {45}
    {30}
\end{problem}

\section{填空题}

\begin{problem}
已知向量 \(\overrightarrow{OA} \perp \overrightarrow{AB}\)，\(|\overrightarrow{OA}| = 3\)，则 \(\overrightarrow{OA} \cdot \overrightarrow{OB} =\)  \fillinblank{}.
\end{problem}

\begin{problem}
若变量 \(x, y\) 满足约束条件 \(\left\{ \begin{array}{l} x + y \leqslant 4, \\ x - y \leqslant 2, \\ 3x - y \geqslant 0, \end{array} \right.\) 则 \(3x + y\) 的最大值是 \fillinblank{}.
\end{problem}

\begin{problem}
函数 \( f(x) = 2\sin x\sin \left(x + \frac{\pi}{2}\right) - x^2 \) 的零点个数为 \fillinblank{}.
\end{problem}

\begin{problem}
某电子商务公司对 \(10000\) 名网络购物者 \(2014\) 年度的消费情况进行统计，发现消费金额（单位：万元）都在区间 \([0.3,0.9]\) 内，其频率分布直方图如图所示. （1）直方图中的 \(a=\fillinblank{}\). （2）在这些购物者中，消费金额在区间 \([0.5,0.9]\) 内的购物者的人数为 \fillinblank{}.
\end{problem}

\begin{problem}
如图，一辆汽车在一条水平公路上向西正行驶，到 \(A\) 处时测得公路北侧一山顶 \(D\) 在西偏北 \(30^{\circ}\) 的方向上，行驶 \(600 \mathrm{~m}\) 后到达 \(B\) 处，测得此山顶在西偏北 \(75^{\circ}\) 的方向上，仰角为 \(30^{\circ}\)，则此山的高度 \(CD = \fillinblank{}\mathrm{m}\).
\end{problem}

\begin{problem}
如图，圆 C 与 x 轴相切于点 T(1,0)，与 y 轴正半轴交于两点 A, B (B 在 A 的上方)，且 \( |AB| = 2 \) .

\((1)\) 圆 C 的标准方程为 \fillinblank{}；\((2)\) 圆 C 在点 B 处的切线在 \(x\) 轴上的截距为 \fillinblank{}.
\end{problem}

\begin{problem}
\(a\) 为实数，函数 \(f(x) = |x^2 - ax|\) 在区间 \([0,1]\) 上的最大值记为 \(g(a)\). 当 \(a = \fillinblank{}\) 时，\(g(a)\) 的值最小.
\end{problem}

\section{解答题}

\begin{problem}
某同学用“五点法”画函数 \( f(x) = A \sin (\omega x + \varphi) \left( \omega > 0, |\varphi| < \frac{\pi}{2} \right) \) 在某一个周期内的图象时，列表并填入了部分数据，如下表： \( \omega x + \varphi \) 0 \( \frac{\pi}{2} \) \( \pi \) \( \frac{3\pi}{2} \) 2π x \( \frac{\pi}{3} \) \( \frac{5\pi}{6} \) \( A\sin(\omega x + \varphi) \) 0 5 -5 0 (1) 请将上表数据补充完整，并直接写出函数 \( f(x) \) 的解析式; (2) 将 \( y = f(x) \) 图象上所有点向左平行移动 \( \frac{\pi}{6} \) 个单位长度，得到 \( y = g(x) \) 的图象. 求 \( y = g(x) \) 的图象离原点 O 最近的对称中心.
\end{problem}

\begin{problem}
设等差数列 \(\{n_n\}\) 的公差为 \(d\)，前 \(n\) 项和为 \(S_n\). 等比数列 \(\{b_n\}\) 的公比为 \(q\). 已知 \(b_1 = a_1, b_2 = 2, q = d, S_{10} = 100\). (1) 求数列 \( \{a_{n}\} \) ，\( \{b_{n}\} \) 的通项公式; (2) 当 \(d > 1\) 时，记 \(c_{n} = \frac{b_{n}}{b_{n}}\)，求数列 \(\{c_{n}\}\) 的前 \(n\) 项和为 \(T_{n}\).
\end{problem}

\begin{problem}
设函数 \( f(x), g(x) \) 的定义域均为 \( \R \)，且 \( f(x) \) 是奇函数，\( g(x) \) 是偶函数，\( f(x) + g(x) = \e^x \)，其中 \( \e \) 为自然对数的底数. (1) 求 \( f(x) \) ，\( g(x) \) 的解析式，并证明: 当 x > 0 时，\( f(x) > 0 \) ，\( g(x) > 1 \) ; (2) 设 \(a \leqslant 0, b \geqslant 1\)，证明：当 \(x > 0\) 时，\(ag(x) + (1 - a) < \frac{f(x)}{x} < bg(x) + (1 - b)\).
\end{problem}

\begin{problem}
一种作图工具如图\circled{1}所示. O 是滑槽 AB 的中点，矩杆 ON 可绕 O 转动，长杆 MN 通过 N 处铰链与 ON 连接，MN 上的栓子 D 可沿滑槽 AB 滑动，且 DN = ON = 1, MN = 3. 当栓子 D 在滑槽 AB 内作往复运动时，带动 N 绕 O 转动一周（D 不动时，N 也不动），M 处的笔尖画出的曲线记为 C. 以 O 为原点，AB 所在的直线为 z 轴建立如图\circled{2}所示的平面直角坐标系. (1) 求曲线 C 的方程; (2) 设动直线 l 与两定直线 \( l_{1}: x - 2y = 0 \) 和 \( l_{2}: x + 2y = 0 \) 分别交于 P, Q 两点. 若直线 l 总与曲线 C 有且只有一个公共点. 试探究: \( \triangle OPQ \) 的面积是否存在最小值\(\perp\) 若存在，求出该最小值; 若不存在，说明理由.

\circled{1}

\circled{2}
\end{problem}

